% GENERATED FILE. Edit canonical lessons, then run npm run generate:books.
% canonical-source-hash: fnv1a64:feefa72693e0f65f
% canonical-lessons: FA-C445-taqato, FA-C445-falake, FA-C445-serah, FA-C445-dorahi, FA-C445-obur-kardan

\chapter{Junctions and roundabouts}
\label{ch:fa-junctions-and-roundabouts}
\hlchaptermodality{445}

\begin{chapteropening}
\textbf{By the end of this chapter:} \emph{I can describe junctions: a crossroads, a roundabout, a T-junction, a fork and crossing over.}

Complete the last lesson of chapter 445: I can describe junctions: a crossroads, a roundabout, a T-junction, a fork and crossing over.
\end{chapteropening}

\section[\texorpdfstring{\fa{تقاطع} (taqâto')}{تقاطع (taqâto')}]{\fa{تقاطع} (taqâto') — a crossroads, an intersection}

\label{lesson:FA-C445-taqato}



\begin{quote}
Before the new one: say the Persian for to get lost, then the Persian for a landmark.
\end{quote}

\subsection*{You'll want to know: \fa{تقاطع}}
\textbf{\fa{تقاطع}} (\emph{taqâto'}) — \textquotedblleft{}a crossroads, an intersection\textquotedblright{}.

\begin{quote}
\fa{تاکسی} \fa{در} \fa{تقاطع} \fa{دوم} \fa{به} \fa{راست} \fa{پیچید}.
\end{quote}

\begin{quote}
\emph{Tâksi dar taqâto'-e dovvom be râst pichid.}
\end{quote}

\begin{quote}
The taxi turned right at the second crossroads.
\end{quote}

More formal than \emph{chahârrâh}.

\begin{cousinweb}
\emph{taqâto'} was borrowed from Arabic \emph{taqāṭuʿ}, which EtymDB lists; it is usually glossed 'cutting across each other'.
\end{cousinweb}

\begin{grammarlens}[title={What you've built}]
One more word for telling, explaining and describing at B1.
\end{grammarlens}

\subsection*{Your turn}
\begin{itemize}
\raggedright
  \item \emph{Say it:} \emph{taqâto'}
  \item \emph{Say it:} \emph{taqâto'}, once more
  \item \emph{Say it:} a sentence of your own with \emph{taqâto'}
  \item \emph{Recall:} say the Persian for to get lost, then the Persian for a landmark, then say \emph{taqâto'} again
\end{itemize}

\begin{tcolorbox}[breakable,colback=teal!4,colframe=teal!35!black,title={Before you move on}]
What does \fa{تقاطع} mean? (\textquotedblleft{}A crossroads, an intersection\textquotedblright{}.) Say it once more.
\end{tcolorbox}
\section[\texorpdfstring{\fa{فلکه} (falake)}{فلکه (falake)}]{\fa{فلکه} (falake) — a roundabout (US: a traffic circle)}

\label{lesson:FA-C445-falake}



\begin{quote}
Before the new one: say the Persian for a landmark, then the Persian for a crossroads.
\end{quote}

\subsection*{You'll want to know: \fa{فلکه}}
\textbf{\fa{فلکه}} (\emph{falake}) — \textquotedblleft{}a roundabout (US: a traffic circle)\textquotedblright{}.

\begin{quote}
\fa{هتل} \fa{ما} \fa{کنار} \fa{فلکه} \fa{اول} \fa{است}.
\end{quote}

\begin{quote}
\emph{Hotel-e mâ kenâr-e falake-ye avval ast.}
\end{quote}

\begin{quote}
Our hotel is next to the first roundabout.
\end{quote}

Many big roundabouts are also called \emph{meydân}.

\begin{cousinweb}
\emph{falake} is usually traced to Arabic \emph{falak} 'the turning sphere of the sky, a circle', with the Persian ending \emph{-e}. EtymDB does not list it.
\end{cousinweb}

\begin{grammarlens}[title={What you've built}]
One more word for telling, explaining and describing at B1.
\end{grammarlens}

\subsection*{Your turn}
\begin{itemize}
\raggedright
  \item \emph{Say it:} \emph{falake}
  \item \emph{Say it:} \emph{falake}, once more
  \item \emph{Say it:} a sentence of your own with \emph{falake}
  \item \emph{Recall:} say the Persian for a landmark, then the Persian for a crossroads, then say \emph{falake} again
\end{itemize}

\begin{tcolorbox}[breakable,colback=teal!4,colframe=teal!35!black,title={Before you move on}]
What does \fa{فلکه} mean? (\textquotedblleft{}A roundabout (US: a traffic circle)\textquotedblright{}.) Say it once more.
\end{tcolorbox}
\section[\texorpdfstring{\fa{سه} \fa{راه} (serâh)}{سه راه (serâh)}]{\fa{سه} \fa{راه} (serâh) — a T-junction, a three-way junction}

\label{lesson:FA-C445-serah}



\begin{quote}
Before the new one: say the Persian for a crossroads, then the Persian for a roundabout.
\end{quote}

\subsection*{You'll want to know: \fa{سه} \fa{راه}}
\textbf{\fa{سه} \fa{راه}} (\emph{serâh}) — \textquotedblleft{}a T-junction, a three-way junction\textquotedblright{}.

\begin{quote}
\fa{به} \fa{سه} \fa{راه} \fa{رسیدیم} \fa{و} \fa{نمیدانستیم} \fa{کدام} \fa{طرف} \fa{برویم}.
\end{quote}

\begin{quote}
\emph{Be serâh residim va nemidânestim kodâm taraf beravim.}
\end{quote}

\begin{quote}
We reached the junction and didn't know which way to go.
\end{quote}

\emph{se} 'three' + \emph{râh} 'road'. A crossroads is \emph{chahârrâh}.

\begin{cousinweb}
\emph{serâh} is a native compound of \emph{se} 'three', usually traced to PIE \emph{tréyes}, and \emph{râh} 'road', from Middle Persian \emph{rāh} (EtymDB).
\end{cousinweb}

\begin{grammarlens}[title={What you've built}]
One more word for telling, explaining and describing at B1.
\end{grammarlens}

\subsection*{Your turn}
\begin{itemize}
\raggedright
  \item \emph{Say it:} \emph{serâh}
  \item \emph{Say it:} \emph{serâh}, once more
  \item \emph{Say it:} a sentence of your own with \emph{serâh}
  \item \emph{Recall:} say the Persian for a crossroads, then the Persian for a roundabout, then say \emph{serâh} again
\end{itemize}

\begin{tcolorbox}[breakable,colback=teal!4,colframe=teal!35!black,title={Before you move on}]
What does \fa{سه} \fa{راه} mean? (\textquotedblleft{}A T-junction, a three-way junction\textquotedblright{}.) Say it once more.
\end{tcolorbox}
\section[\texorpdfstring{\fa{دوراهی} (dorâhi)}{دوراهی (dorâhi)}]{\fa{دوراهی} (dorâhi) — a fork (in the road); a dilemma}

\label{lesson:FA-C445-dorahi}



\begin{quote}
Before the new one: say the Persian for a roundabout, then the Persian for a T-junction.
\end{quote}

\subsection*{You'll want to know: \fa{دوراهی}}
\textbf{\fa{دوراهی}} (\emph{dorâhi}) — \textquotedblleft{}a fork (in the road); a dilemma\textquotedblright{}.

\begin{quote}
\fa{در} \fa{دوراهی} \fa{راه} \fa{سمت} \fa{چپ} \fa{را} \fa{گرفتیم}.
\end{quote}

\begin{quote}
\emph{Dar dorâhi râh-e samt-e chap râ gereftim.}
\end{quote}

\begin{quote}
At the fork we took the road on the left.
\end{quote}

Also figurative: \emph{sar-e dorâhi budan} — to be at a crossroads in life.

\begin{cousinweb}
\emph{dorâhi} is a native compound of \emph{do} 'two', which EtymDB traces to PIE \emph{dwóh\textsubscript{1}}, \emph{râh} 'road', from Middle Persian \emph{rāh} (EtymDB), and the ending \emph{-i}.
\end{cousinweb}

\begin{grammarlens}[title={What you've built}]
One more word for telling, explaining and describing at B1.
\end{grammarlens}

\subsection*{Your turn}
\begin{itemize}
\raggedright
  \item \emph{Say it:} \emph{dorâhi}
  \item \emph{Say it:} \emph{dorâhi}, once more
  \item \emph{Say it:} a sentence of your own with \emph{dorâhi}
  \item \emph{Recall:} say the Persian for a roundabout, then the Persian for a T-junction, then say \emph{dorâhi} again
\end{itemize}

\begin{tcolorbox}[breakable,colback=teal!4,colframe=teal!35!black,title={Before you move on}]
What does \fa{دوراهی} mean? (\textquotedblleft{}A fork (in the road); a dilemma\textquotedblright{}.) Say it once more.
\end{tcolorbox}
\section[\texorpdfstring{\fa{عبور} \fa{کردن} ('obur kardan)}{عبور کردن ('obur kardan)}]{\fa{عبور} \fa{کردن} ('obur kardan) — to cross, to go across}

\label{lesson:FA-C445-obur-kardan}



\begin{quote}
Before the new one: say the Persian for a T-junction, then the Persian for a fork.
\end{quote}

\subsection*{You'll want to know: \fa{عبور} \fa{کردن}}
\textbf{\fa{عبور} \fa{کردن}} (\emph{'obur kardan}) — \textquotedblleft{}to cross, to go across\textquotedblright{}.

\begin{quote}
\fa{بچه} \fa{ها} \fa{با} \fa{دقت} \fa{از} \fa{خیابان} \fa{عبور} \fa{کردند}.
\end{quote}

\begin{quote}
\emph{Bachche-hâ bâ deqqat az khiyâbân 'obur kardand.}
\end{quote}

\begin{quote}
The children crossed the street carefully.
\end{quote}

More formal than \emph{rad shodan}. \emph{'Obur mamnu'} — no crossing.

\begin{cousinweb}
\emph{'obur kardan} is built from \emph{'obur} 'crossing', a loan usually traced to Arabic \emph{ʿubūr}, and \emph{kardan} 'to do', inherited from Middle Persian \emph{kardan} (EtymDB).
\end{cousinweb}

\begin{grammarlens}[title={What you've built}]
One more word for telling, explaining and describing at B1.
\end{grammarlens}

\subsection*{Your turn}
\begin{itemize}
\raggedright
  \item \emph{Say it:} \emph{'obur kardan}
  \item \emph{Say it:} \emph{'obur kardan}, once more
  \item \emph{Say it:} a sentence of your own with \emph{'obur kardan}
  \item \emph{Recall:} say the Persian for a T-junction, then the Persian for a fork, then say \emph{'obur kardan} again
\end{itemize}

\begin{tcolorbox}[breakable,colback=teal!4,colframe=teal!35!black,title={Before you move on}]
What does \fa{عبور} \fa{کردن} mean? (\textquotedblleft{}To cross, to go across\textquotedblright{}.) Say it once more.
\end{tcolorbox}
